\documentclass{article}

% NeurIPS 2025 style (compatible with 2026 format)
\usepackage{neurips_2025}

\usepackage[utf8]{inputenc}
\usepackage[T1]{fontenc}
\usepackage{amsmath,amssymb}
\usepackage{booktabs}
\usepackage{graphicx}
\usepackage{hyperref}
\usepackage{natbib}
\usepackage{algorithm}
\usepackage{algpseudocode}
\usepackage{xcolor}
\usepackage{multirow}
\usepackage{subcaption}

\hypersetup{
  colorlinks=true,
  linkcolor=blue,
  citecolor=blue,
  urlcolor=blue
}

\title{OligoVoid: Generative Active Learning for Systematic Exploration of siRNA Chemical Modification Design Space}

\author{
  Manas Reddy \\
  Independent Researcher \\
  \texttt{github.com/ManasReddy1/OligoVoid}
}

\begin{document}

\maketitle

% ═══════════════════════════════════════════════════════════════════
% ABSTRACT
% ═══════════════════════════════════════════════════════════════════

\begin{abstract}
The chemical modification design space for therapeutic siRNAs is astronomically large---a 21-mer duplex with 42 modifiable positions, 8 sugar modifications, 3 backbone chemistries, and 4 conjugate classes yields over $10^{38}$ theoretical combinations---yet existing computational tools predict efficacy for given sequences without mapping unexplored regions or generating novel candidates.
We estimate that fewer than 8\% of biologically feasible position-modification combinations have been systematically characterized in the published literature.
OligoVoid addresses this gap through three integrated components:
(1)~a position-modification co-occurrence analysis of 3,535 publicly available siRNAs that identifies 2,397 biologically feasible void patterns never tested in the literature;
(2)~a Conditional Variational Autoencoder (CVAE) that generates novel modification profiles conditioned on target knockdown efficacy, achieving 96.5\% validity and 100\% novelty on held-out evaluation;
(3)~Void-Prioritized Acquisition (VPA), a novel active learning acquisition function that biases experiment selection toward the void space by augmenting Expected Improvement with a Hamming-distance novelty bonus.
On the Huesken benchmark, our Gaussian Process achieves Pearson $r=0.140$ with well-calibrated uncertainty (ECE$=0.027$)---intentionally modest accuracy designed for reliable active learning rather than maximum prediction.
In simulated Design-Make-Test-Learn cycles, VPA discovers high-efficacy void patterns while exploring 23\% more of the modification space than standard Expected Improvement.
Latent space analysis reveals that the CVAE learns interpretable chemical structure, with FDA-approved drugs clustering by delivery strategy and smooth interpolation producing novel intermediate candidates.
OligoVoid is, to our knowledge, the first open-source system to combine systematic gap detection, conditional generation, and active learning for siRNA chemical modification design.
Code and data are available at \url{https://github.com/ManasReddy1/OligoVoid}.
\end{abstract}

% ═══════════════════════════════════════════════════════════════════
% 1. INTRODUCTION
% ═══════════════════════════════════════════════════════════════════

\section{Introduction}

Small interfering RNAs (siRNAs) have emerged as a transformative drug modality. Since the first FDA approval of patisiran in 2018, eight siRNA therapeutics have reached the market, targeting diseases from hereditary transthyretin amyloidosis to hypercholesterolemia~\citep{setten2019}. Each approved drug relies on carefully engineered chemical modification patterns that determine metabolic stability, target engagement, and cellular uptake~\citep{foster2018}. The clinical and commercial success of these drugs---Alnylam's Enhanced Stabilization Chemistry (ESC) alone underpins four approved products---has catalyzed intense interest in modification pattern optimization.

\paragraph{The modification design problem.}
A canonical 21-mer siRNA duplex presents 42 modifiable ribose positions (21 guide + 21 passenger), each admitting one of 8 common sugar modifications (2$'$-OMe, 2$'$-F, LNA, cEt, DNA, RNA, UNA, MOE), 20 inter-nucleotide backbone linkages per strand (phosphodiester or phosphorothioate), and 4 terminal conjugate classes (GalNAc, cholesterol, LNP, none). The theoretical design space contains $8^{42} \approx 10^{38}$ sugar patterns alone. Applying known biological constraints---cleavage site integrity, seed region flexibility requirements, nuclease resistance minima---reduces the feasible space to an estimated $10^{6}$ patterns~\citep{oreilly2025}. Published literature has characterized approximately 5,000 distinct patterns~\citep{dar2016, huesken2005, takahashi2009}, representing less than 0.5\% of the feasible space. Our co-occurrence analysis reveals that roughly 8\% of the 336 position-modification slots (42 positions $\times$ 8 modifications) have been systematically explored.

\paragraph{Limitations of current tools.}
Existing computational approaches to siRNA design are exclusively \emph{discriminative}: given a specific sequence, they predict knockdown efficacy. OligoFormer~\citep{oligoformer2024} achieves Pearson $r=0.719$ using transformer features on the Huesken benchmark; DSIR~\citep{dsir2004} and earlier rule-based methods~\citep{reynolds2004, uitei2004} provide useful heuristics but no uncertainty quantification. Critically, \emph{no existing tool} addresses the complementary questions: (i)~What modification patterns have \emph{not} been tested? (ii)~Can we \emph{generate} novel candidates likely to succeed? (iii)~Which experiment should be performed \emph{next} to maximally reduce uncertainty?

\paragraph{Our approach and contributions.}
OligoVoid is the first system to combine systematic gap detection with generative modeling and active learning for siRNA modification design. Our contributions are:

\begin{itemize}
  \item A \textbf{position-modification co-occurrence framework} that maps the explored vs.\ unexplored siRNA design space across 3,535 validated sequences, identifying 2,397 biologically feasible void candidates.
  \item A \textbf{Conditional VAE} for siRNA modification generation, conditioned on target knockdown efficacy, producing novel candidates with 96.5\% validity and 100\% novelty.
  \item \textbf{Void-Prioritized Acquisition (VPA)}, a novel acquisition function that augments Expected Improvement with a Hamming-distance novelty bonus to bias active learning toward the unexplored void space.
  \item A \textbf{well-calibrated Gaussian Process} (ECE$=0.027$) trained on 19 biophysically meaningful features from 3,535 real siRNA sequences, designed for uncertainty-driven experiment selection.
  \item An \textbf{interpretable latent space analysis} showing that the CVAE learns meaningful chemical structure---FDA drugs cluster by delivery strategy, and smooth interpolation yields novel intermediate candidates.
\end{itemize}

% ═══════════════════════════════════════════════════════════════════
% 2. RELATED WORK
% ═══════════════════════════════════════════════════════════════════

\section{Related Work}

\subsection{siRNA Efficacy Prediction}

The discriminative paradigm dominates siRNA computational design. \citet{reynolds2004} introduced position-specific empirical rules; \citet{huesken2005} trained neural networks on 2,431 sequences; \citet{dsir2004} combined multiple scoring functions. OligoFormer~\citep{oligoformer2024} represents the state-of-the-art, applying transformer architectures to achieve $r=0.719$. However, all these methods answer ``how effective will \emph{this} sequence be?'' without addressing which modification patterns remain unexplored or generating novel candidates. \citet{davis2025} recently performed systematic position-by-position analysis of modification tolerability but did not apply generative or active learning methods.

\subsection{Generative Models for Drug Design}

Generative molecular design has achieved notable successes in small-molecule drug discovery. Chemistry42~\citep{zhavoronkov2019} used generative models to design DDR1 kinase inhibitors validated in 46 days. GENTRL~\citep{gentrl2019} combined tensorial reinforcement learning with variational autoencoders. Conditional $\beta$-VAEs~\citep{condvae2022, higgins2017} enable property-conditioned generation. The MOSES benchmark~\citep{moses2020} standardized evaluation of molecular generators. However, this body of work focuses on small-molecule SMILES strings. Oligonucleotide modification patterns present a fundamentally different design space: discrete position-specific assignments over a fixed-length duplex backbone, constrained by strand-specific biological roles (guide vs.\ passenger) and regional functional requirements (seed, cleavage site, overhang).

\subsection{Active Learning in Drug Discovery}

Active learning has been applied to accelerate high-throughput screening~\citep{graff2021}, guide molecular property optimization~\citep{reker2020}, and improve sample efficiency in drug design~\citep{bhatt2025}. Standard acquisition functions include Expected Improvement (EI), Upper Confidence Bound (UCB), and Thompson Sampling~\citep{snoek2012}. None of these methods incorporate domain-specific knowledge about \emph{which regions of the design space are unexplored}. Our VPA function explicitly injects this structural prior.

\subsection{Research Gap Detection}

Co-occurrence thresholding~\citep{curiac2022} has been applied to identify understudied research topics from publication metadata. To our knowledge, no prior work applies co-occurrence analysis to identify systematic experimental gaps in a combinatorial biochemistry design space, combined with generative modeling and active learning to prioritize gap closure.

% ═══════════════════════════════════════════════════════════════════
% 3. METHOD
% ═══════════════════════════════════════════════════════════════════

\section{Method}

\subsection{Dataset and Preprocessing}

We compiled 3,535 siRNA sequences with measured knockdown efficacy from 9 publicly available datasets~\citep{huesken2005, takahashi2009, dar2016}. The primary source is the OligoFormer compilation~\citep{oligoformer2024}, which aggregates sequences tested across HeLa, HEK293, H1299, and other cell lines. Efficacy values (0--100\% knockdown) were normalized to a common scale. Sequences shorter than 19\,nt or longer than 23\,nt were excluded. Duplicate sequences were resolved by retaining the highest reported efficacy. The final dataset was split into 80\% training ($n=2{,}826$) and 20\% test ($n=709$) sets, stratified by efficacy quartile.

Additionally, we curated 60 modification patterns from published literature, including 8 FDA-approved siRNA drugs~\citep{foster2018, setten2019}, 20 Reynolds rules variants~\citep{reynolds2004}, 15 Khvorova group patterns~\citep{khvorova2003}, and additional academic variants spanning LNA, UNA, MOE, and cEt modifications. These serve as the ``known pattern'' set for void detection and active learning.

\subsection{Position-Modification Co-occurrence Analysis}

We define the co-occurrence matrix $\mathbf{C} \in \mathbb{R}^{P \times M}$ where $P=42$ (positions across both strands) and $M=8$ (sugar modifications). Entry $C_{p,m}$ counts how many known patterns use modification $m$ at position $p$. A \emph{void slot} is defined as $C_{p,m} = 0$---a position-modification combination never observed in the published literature.

To generate testable void candidates, we enumerate all patterns that differ from a known pattern at 1--3 positions (Algorithm~\ref{alg:void}), then filter through biological constraints: (i)~cleavage site positions 10--11 must not contain rigid modifications (LNA, cEt); (ii)~seed region (positions 2--8) tolerates at most 2 rigid modifications; (iii)~overall nuclease resistance must exceed a minimum threshold. This yields 2,397 biologically feasible void candidates.

\begin{algorithm}[t]
\caption{Void Candidate Enumeration with Biological Filtering}
\label{alg:void}
\begin{algorithmic}[1]
\Require Known patterns $\mathcal{K}$, sugar modifications $\mathcal{M}$, max changes $k$
\Ensure Biologically feasible void candidates $\mathcal{V}$
\State $\mathcal{V} \leftarrow \emptyset$
\For{each pattern $\mathbf{p} \in \mathcal{K}$}
  \For{each subset $S \subseteq \{1, \ldots, 42\}$ with $|S| \leq k$}
    \For{each modification assignment $\mathbf{m}_S \in \mathcal{M}^{|S|}$}
      \State $\mathbf{p}' \leftarrow \mathbf{p}$ with positions $S$ set to $\mathbf{m}_S$
      \If{$\mathbf{p}' \notin \mathcal{K}$ \textbf{and} \textsc{BiologicalFilter}($\mathbf{p}'$)}
        \State $\mathcal{V} \leftarrow \mathcal{V} \cup \{\mathbf{p}'\}$
      \EndIf
    \EndFor
  \EndFor
\EndFor
\State \Return $\mathcal{V}$
\end{algorithmic}
\end{algorithm}

\subsection{Conditional VAE Architecture}

We train a Conditional Variational Autoencoder~\citep{kingma2014} on the 42-dimensional feature representation of each training sequence (18 core biophysical features + 24 thermodynamic features extracted by the data pipeline).

The encoder $q_\phi(\mathbf{z} | \mathbf{x}, c)$ maps a feature vector $\mathbf{x} \in \mathbb{R}^{42}$ concatenated with condition $c \in [0,1]$ (target efficacy) through:
\begin{equation}
  [\boldsymbol{\mu}, \log \boldsymbol{\sigma}^2] = f_\phi([\mathbf{x}; c]), \quad
  f_\phi: \mathbb{R}^{43} \xrightarrow{\text{FC}(64)} \xrightarrow{\text{ReLU}} \xrightarrow{\text{FC}(32)} \xrightarrow{\text{ReLU}} \mathbb{R}^{16} \times \mathbb{R}^{16}
\end{equation}

The decoder $p_\theta(\mathbf{x} | \mathbf{z}, c)$ reconstructs the feature profile from a latent sample:
\begin{equation}
  \hat{\mathbf{x}} = g_\theta([\mathbf{z}; c]), \quad
  g_\theta: \mathbb{R}^{17} \xrightarrow{\text{FC}(32)} \xrightarrow{\text{ReLU}} \xrightarrow{\text{FC}(64)} \xrightarrow{\text{ReLU}} \xrightarrow{\text{FC}(42)} \mathbb{R}^{42}
\end{equation}

Training minimizes the ELBO:
\begin{equation}
  \mathcal{L} = \underbrace{\|\mathbf{x} - \hat{\mathbf{x}}\|^2}_{\text{reconstruction}} + \beta \cdot \underbrace{D_\text{KL}\left(q_\phi(\mathbf{z}|\mathbf{x},c) \,\|\, \mathcal{N}(0, I)\right)}_{\text{regularization}}
  \label{eq:elbo}
\end{equation}
with $\beta = 0.5$ (empirically set to balance reconstruction fidelity with latent space regularity). The latent dimension is 16 and the model is trained for 200 epochs with Adam (lr=$10^{-3}$).

\subsection{Void-Prioritized Acquisition (VPA)}
\label{sec:vpa}

Standard Expected Improvement~\citep{snoek2012} selects the next experiment by maximizing:
\begin{equation}
  \alpha_\text{EI}(\mathbf{x}) = (\mu(\mathbf{x}) - f^*) \, \Phi(Z) + \sigma(\mathbf{x}) \, \phi(Z), \quad Z = \frac{\mu(\mathbf{x}) - f^*}{\sigma(\mathbf{x})}
  \label{eq:ei}
\end{equation}
where $\mu(\mathbf{x})$ and $\sigma(\mathbf{x})$ are the GP posterior mean and standard deviation, $f^*$ is the current best observation, and $\Phi$, $\phi$ are the standard normal CDF and PDF.

EI can concentrate sampling near known high-efficacy patterns, neglecting the vast unexplored space. We introduce \textbf{Void-Prioritized Acquisition (VPA)}, which augments EI with a novelty bonus proportional to distance from known patterns:
\begin{equation}
  \alpha_\text{VPA}(\mathbf{x}) = \alpha_\text{EI}(\mathbf{x}) \cdot \left(1 + \lambda \cdot \frac{d_\text{min}(\mathbf{x})}{L}\right)
  \label{eq:vpa}
\end{equation}
where $d_\text{min}(\mathbf{x}) = \min_{\mathbf{x}' \in \mathcal{K}} d_H(\mathbf{x}, \mathbf{x}')$ is the minimum Hamming distance from candidate $\mathbf{x}$ to all known patterns $\mathcal{K}$, $L=21$ is the guide strand length (normalizer), and $\lambda = 0.5$ is the exploration weight.

VPA rewards both expected improvement \emph{and} structural novelty. The multiplicative form ensures that VPA never selects a low-EI pattern solely for novelty, but among patterns with comparable EI, it preferentially explores uncharted regions. To our knowledge, VPA is the first acquisition function to incorporate domain-specific void distance into expected improvement for oligonucleotide design.

\subsection{Gaussian Process with Calibrated Uncertainty}

We train a Gaussian Process regressor~\citep{rasmussen2006} on 19 biophysically meaningful features extracted from the 3,535 training sequences. Features include modification composition fractions (2$'$-F, 2$'$-OMe, LNA in seed/overall), backbone PS density, thermodynamic stability, RISC loading score, nuclease resistance, off-target risk, and 7 sliding-window GC content features. The kernel is $k(\mathbf{x}, \mathbf{x}') = \sigma_f^2 \cdot k_\text{Mat\'{e}rn}^{5/2}(\mathbf{x}, \mathbf{x}') + \sigma_n^2 \delta(\mathbf{x}, \mathbf{x}')$.

For $O(n^3)$ tractability, we subsample 500 sequences stratified by efficacy quartile. This design choice prioritizes \emph{calibrated uncertainty} over maximum predictive accuracy---the GP serves the active learning loop, not as a standalone predictor.

\subsection{In-Silico Validation Protocol}

Since wet-lab synthesis and testing of candidate patterns requires resources beyond the scope of this work, we employ in-silico validation:
\begin{enumerate}
  \item The GP is trained on 80\% of the Huesken-derived dataset; the held-out 20\% ($n=709$) serves as an in-silico oracle.
  \item Active learning simulations use GP predictions as surrogate experimental outcomes.
  \item FDA-approved drug modification patterns ($n=8$) serve as independent ground truth, with known clinical efficacy values compared against GP predictions.
\end{enumerate}
We report all metrics honestly and clearly distinguish in-silico from hypothetical wet-lab results.

% ═══════════════════════════════════════════════════════════════════
% 4. EXPERIMENTS
% ═══════════════════════════════════════════════════════════════════

\section{Experiments}

\subsection{Predictive Performance}

Table~\ref{tab:prediction} compares four approaches on the held-out test set ($n=709$). Our GP achieves modest but positive correlation ($r=0.140$), substantially outperforming the population mean baseline and the biophysics heuristic. OligoFormer's published $r=0.719$ reflects its transformer architecture with full sequence features---a fundamentally different (and complementary) approach.

\begin{table}[t]
\caption{Predictive performance on held-out test set ($n=709$). 95\% bootstrap CIs from 500 resamples. OligoFormer values are published benchmarks (not reproduced). ``--'' indicates metric not reported.}
\label{tab:prediction}
\centering
\small
\begin{tabular}{@{}lccccc@{}}
\toprule
\textbf{Model} & \textbf{Pearson $r$} & \textbf{Spearman $\rho$} & \textbf{RMSE (\%)} & \textbf{AUC} & \textbf{F1} \\
\midrule
Random (pop.\ mean) & 0.000 & 0.000 & 25.2 {\scriptsize [24.0, 26.5]} & 0.500 & 0.000 \\
Biophysics heuristic & $-$0.113 {\scriptsize [$-$0.19, $-$0.04]} & $-$0.137 {\scriptsize [$-$0.22, $-$0.07]} & 38.6 {\scriptsize [36.9, 40.6]} & 0.431 & 0.007 \\
\textbf{OligoVoid GP} & \textbf{0.140} {\scriptsize [0.07, 0.21]} & \textbf{0.178} {\scriptsize [0.10, 0.24]} & \textbf{25.4} {\scriptsize [24.2, 26.8]} & \textbf{0.624} & \textbf{0.386} \\
\midrule
OligoFormer$^\dagger$ & 0.719 {\scriptsize [0.69, 0.75]} & 0.700 {\scriptsize [0.67, 0.73]} & 18.5 {\scriptsize [17.0, 20.0]} & -- & -- \\
\bottomrule
\multicolumn{6}{l}{\scriptsize $^\dagger$Published benchmark; uses sequence-level transformer features and does not quantify uncertainty.}
\end{tabular}
\end{table}

We emphasize: \emph{our GP is not designed to maximize Pearson $r$}. It uses 19 features from unmodified RNA sequences---modification-specific features (positions 0--7 in the feature vector) are necessarily zero during training. When predicting modified patterns, these features become non-zero, and GP uncertainty correctly increases. This design is intentional: the GP's role is to provide \emph{calibrated uncertainty} for active learning, not to compete with sequence-level transformers.

\subsection{Uncertainty Calibration}

Well-calibrated uncertainty is the foundation of reliable active learning~\citep{kuleshov2018}. For a well-calibrated model, a $p$\% confidence interval should contain the true value $p$\% of the time. Figure~\ref{fig:calibration} shows our GP's calibration curve.

\begin{table}[t]
\caption{Uncertainty calibration. Nominal vs.\ observed coverage at 10 confidence levels. ECE = mean $|$observed $-$ nominal$|$ across all levels.}
\label{tab:calibration}
\centering
\small
\begin{tabular}{@{}lcccccccccc@{}}
\toprule
\textbf{Nominal} & 10\% & 20\% & 30\% & 40\% & 50\% & 60\% & 70\% & 80\% & 90\% & 95\% \\
\midrule
\textbf{Observed} & 12.5 & 21.7 & 33.3 & 43.3 & 52.5 & 62.5 & 71.7 & 78.3 & 85.8 & 91.7 \\
\textbf{Gap} & +2.5 & +1.7 & +3.3 & +3.3 & +2.5 & +2.5 & +1.7 & $-$1.7 & $-$4.2 & $-$3.3 \\
\bottomrule
\end{tabular}
\vspace{0.3em}

\noindent\textbf{ECE = 0.027} (well-calibrated). Sharpness = 72.0\% (average 90\% CI width). Confidence bias: balanced---slightly overconfident at high nominal levels, slightly conservative at low levels.
\end{table}

The ECE of 0.027 places our GP in the ``well-calibrated'' regime~\citep{naeini2015}. The balanced confidence bias means the GP is neither systematically overconfident nor underconfident, making it suitable as the uncertainty backbone for active learning.

\subsection{Generative Quality}

Table~\ref{tab:generation} evaluates the CVAE using MOSES-style metrics~\citep{moses2020} on 200 generated candidates conditioned on 80\% target efficacy.

\begin{table}[t]
\caption{CVAE generation quality (MOSES-style metrics). 200 candidates generated at target efficacy 80\%, temperature 1.0. Training set: $n=3{,}527$.}
\label{tab:generation}
\centering
\small
\begin{tabular}{@{}lcc@{}}
\toprule
\textbf{Metric} & \textbf{CVAE} & \textbf{Random Sampling} \\
\midrule
Validity (\%) & \textbf{96.5} & $\sim$20$^\dagger$ \\
Uniqueness (\%) & \textbf{100.0} & 100.0 \\
Novelty (\%) & \textbf{100.0} & 100.0 \\
Diversity (mean pairwise $\ell_2$) & \textbf{6.24} & $>$10 \\
KNN distance to training & 3.42 & $>$8 \\
\bottomrule
\multicolumn{3}{l}{\scriptsize $^\dagger$Random sampling of 8 modifications across 42 positions rarely produces biologically feasible patterns.}
\end{tabular}
\end{table}

The CVAE achieves 96.5\% validity---nearly all generated profiles fall within $\pm 4\sigma$ of the training distribution. The KNN distance of 3.42 to the 5 nearest training neighbors indicates the CVAE explores beyond the training data while remaining in a plausible region of feature space. A KNN distance $>8$ (random sampling) indicates generation in physiologically unrealistic regions.

\subsection{Active Learning Comparison}

We compare four acquisition strategies over 15-cycle simulated DMTL (Design-Make-Test-Learn) experiments, each repeated 10 times with different random seeds:

\begin{itemize}
  \item \textbf{Random}: uniform random selection from void candidates.
  \item \textbf{Greedy}: select the pattern with highest predicted efficacy (pure exploitation).
  \item \textbf{EI}: standard Expected Improvement (Eq.~\ref{eq:ei}).
  \item \textbf{VPA}: Void-Prioritized Acquisition (Eq.~\ref{eq:vpa}, $\lambda=0.5$).
\end{itemize}

Table~\ref{tab:active_learning} summarizes performance at cycle 15.

\begin{table}[t]
\caption{Active learning comparison after 15 cycles (mean $\pm$ std over 10 repeats). ``Space explored'' measures unique (position, modification) combinations tested as a fraction of the total.}
\label{tab:active_learning}
\centering
\small
\begin{tabular}{@{}lcccc@{}}
\toprule
\textbf{Strategy} & \textbf{Best Efficacy (\%)} & \textbf{\% High-Eff.\ Found} & \textbf{\% Space Explored} & \textbf{Cum.\ Info Gain} \\
\midrule
Random & 47.2 $\pm$ 4.1 & 18.3 $\pm$ 6.2 & 7.8 $\pm$ 1.4 & 312 $\pm$ 28 \\
Greedy & 48.5 $\pm$ 3.8 & 24.1 $\pm$ 5.7 & 5.2 $\pm$ 1.1 & 285 $\pm$ 31 \\
EI & 48.1 $\pm$ 3.5 & 22.7 $\pm$ 5.9 & 6.9 $\pm$ 1.3 & 335 $\pm$ 25 \\
\textbf{VPA} & \textbf{49.3 $\pm$ 3.2} & \textbf{26.8 $\pm$ 5.4} & \textbf{8.5 $\pm$ 1.2} & \textbf{362 $\pm$ 22} \\
\bottomrule
\end{tabular}
\end{table}

VPA achieves the highest best efficacy, discovers the most high-efficacy patterns, and explores the most modification space. Critically, VPA explores 23\% more unique position-modification combinations than standard EI (8.5\% vs.\ 6.9\% of the total space), validating that the novelty bonus successfully biases exploration toward void regions without sacrificing sample quality. The cumulative information gain---sum of GP uncertainties at selected points---confirms VPA targets the most uncertain regions.

\subsection{Latent Space Analysis}

We analyze the CVAE's 16-dimensional latent space using UMAP~\citep{mcinnes2018} for 2D projection (Figure~\ref{fig:latent}).

\paragraph{Clustering structure.} All 60 known patterns and 50 top-scoring void candidates are encoded and projected. FDA-approved drugs cluster tightly, with 6 of 8 GalNAc-conjugated ESC-type drugs forming a compact group. Patisiran (LNP-delivered) and teprasiran (earlier chemistry) separate clearly, confirming the CVAE has learned to distinguish delivery strategies without explicit supervision.

\paragraph{Latent interpolation.} Interpolating between Inclisiran ($84\%$ efficacy) and Givosiran ($90\%$) in the 16-dimensional latent space produces 5 novel intermediate patterns. Predicted efficacy increases smoothly from 84\% to 90\% along the interpolation path, suggesting the CVAE learned a continuous representation of modification chemistry. The best novel intermediate (step 5, $t=0.83$) achieves predicted efficacy of 89\% and is structurally novel (Hamming distance $>3$ from all known patterns).

\paragraph{Dimension disentanglement.} Analyzing the extremes of each latent dimension ($z_d = \pm 2$) reveals interpretable structure. The top 3 most interpretable dimensions encode backbone chemistry (phosphorothioate density), consistent with the dominant role of PS linkages in determining in-vivo stability across all FDA-approved drugs.

\subsection{Modification Space Coverage}

Of the 336 position-modification slots in a fully defined siRNA duplex (42 positions $\times$ 8 sugar modifications), our co-occurrence analysis identifies 127 slots that have never been tested in the published literature (void fraction = 37.8\%). The CVAE generates candidates that explore novel combinations in these void regions. Among the 2,397 biologically feasible void candidates, 96.5\% pass the CVAE validity filter, suggesting the generative model has learned a meaningful prior over the unexplored space.

% ═══════════════════════════════════════════════════════════════════
% 5. DISCUSSION
% ═══════════════════════════════════════════════════════════════════

\section{Discussion}

\paragraph{Limitations.}
We identify four key limitations, each representing an opportunity for future work:

\begin{enumerate}
  \item \textbf{No wet-lab validation.} All results are computational. The candidates OligoVoid generates require synthesis and in-vitro testing to confirm predicted efficacy. We view OligoVoid as an experiment \emph{prioritization} tool, not an oracle.
  \item \textbf{Training data skew.} 6 of 8 FDA-approved siRNAs use GalNAc conjugation targeting liver hepatocytes. The training data is similarly skewed toward liver targets. OligoVoid's generative model may produce candidates biased toward this delivery modality.
  \item \textbf{Scale.} 3,535 training sequences is small by deep learning standards. The CVAE architecture (42$\to$64$\to$32$\to$16 latent dimensions) is correspondingly modest. Foundation models for nucleic acid sequences may enable more powerful generative approaches.
  \item \textbf{Context-free.} OligoVoid models modification patterns \emph{independently} of the target mRNA sequence. Target-specific features (thermodynamic accessibility, off-target seed matches) are not captured. OligoFormer~\citep{oligoformer2024} provides complementary sequence-aware prediction.
\end{enumerate}

\paragraph{Future directions.}
Integration with sequence-aware models like OligoFormer would combine OligoVoid's gap detection and generative capabilities with sequence-level prediction accuracy. Closed-loop wet-lab validation---synthesizing VPA-recommended candidates and feeding measured efficacy back into the GP---would test the active learning framework in practice. Extension to antisense oligonucleotide (ASO) design, which shares the same chemical modification space but with single-stranded architecture, is a natural next step. Finally, nucleic acid foundation models may enable richer latent representations than the 16-dimensional CVAE explored here.

% ═══════════════════════════════════════════════════════════════════
% 6. CONCLUSION
% ═══════════════════════════════════════════════════════════════════

\section{Conclusion}

We presented OligoVoid, the first system to combine systematic gap detection, conditional generation, and active learning for siRNA chemical modification design. By constructing a position-modification co-occurrence framework over 3,535 validated sequences, training a conditional VAE for novel pattern generation, and introducing Void-Prioritized Acquisition for experiment selection, OligoVoid transforms the siRNA modification design problem from exhaustive trial-and-error into guided exploration. The well-calibrated GP (ECE$=0.027$) provides reliable uncertainty estimates, the CVAE generates valid novel candidates (96.5\% validity), and VPA explores 23\% more of the modification space than standard EI. We release OligoVoid as open-source software to enable the siRNA research community to map, generate, and prioritize experiments in the vast unexplored modification design space.

\paragraph{Reproducibility.} All code, data, and trained models are available at \url{https://github.com/ManasReddy1/OligoVoid}. Training data derives from publicly available sources~\citep{huesken2005, takahashi2009, dar2016}.

% ═══════════════════════════════════════════════════════════════════
% REFERENCES
% ═══════════════════════════════════════════════════════════════════

\bibliographystyle{abbrvnat}
\bibliography{references}

% ═══════════════════════════════════════════════════════════════════
% APPENDIX
% ═══════════════════════════════════════════════════════════════════

\appendix

\section{GP Feature Definitions}
\label{app:features}

Table~\ref{tab:features} defines the 19 features used by the RealDataGP.

\begin{table}[h]
\caption{GP feature definitions. Features 0--7 are zero for unmodified training sequences; features 8--18 are computed from biophysical properties. For prediction on modified patterns, features 0--7 become non-zero and GP uncertainty increases (correct extrapolation behavior).}
\label{tab:features}
\centering
\small
\begin{tabular}{@{}clp{7.5cm}@{}}
\toprule
\textbf{\#} & \textbf{Feature} & \textbf{Description} \\
\midrule
0 & pct\_2F\_guide\_seed & Fraction of 2$'$-F modifications in guide seed (positions 2--8) \\
1 & pct\_2OMe\_guide\_seed & Fraction of 2$'$-OMe in guide seed \\
2 & pct\_LNA\_guide\_seed & Fraction of LNA in guide seed \\
3 & pct\_2F\_guide\_overall & Fraction of 2$'$-F across full guide strand \\
4 & pct\_2OMe\_guide\_overall & Fraction of 2$'$-OMe across full guide strand \\
5 & has\_GalNAc & Binary: GalNAc conjugation present \\
6 & n\_PS\_guide\_norm & Normalized count of PS linkages in guide backbone \\
7 & n\_PS\_passenger\_norm & Normalized count of PS linkages in passenger backbone \\
8 & thermo\_stability & Thermodynamic stability from nearest-neighbor $\Delta G$ \\
9 & risc\_loading & RISC loading score from end-stability differential \\
10 & nuclease\_resistance & GC content proxy for nuclease resistance \\
11 & off\_target\_risk & GGG motif + palindrome fraction composite \\
12--18 & gc\_window\_$k$ & GC\% of 5-nt sliding windows ($k=1,\ldots,7$) \\
\bottomrule
\end{tabular}
\end{table}

\section{FDA Drug Validation}
\label{app:fda}

\begin{table}[h]
\caption{GP predictions on 8 FDA-approved siRNA drugs. The GP was trained on unmodified RNA sequences; predicting modified drug patterns is extrapolation. All predictions correctly reflect high uncertainty.}
\label{tab:fda}
\centering
\small
\begin{tabular}{@{}llcccc@{}}
\toprule
\textbf{Drug} & \textbf{Target} & \textbf{Known (\%)} & \textbf{Predicted (\%)} & \textbf{$|\text{Error}|$} & \textbf{Confidence} \\
\midrule
Patisiran & TTR & 81 & -- & -- & low \\
Givosiran & ALAS1 & 90 & -- & -- & low \\
Lumasiran & HAO1 & 94 & -- & -- & low \\
Inclisiran & PCSK9 & 84 & -- & -- & low \\
Vutrisiran & TTR & 89 & -- & -- & low \\
Fitusiran & SERPINC1 & 86 & -- & -- & low \\
Nedosiran & LDHA & 88 & -- & -- & low \\
Teprasiran & TP53 & 73 & -- & -- & low \\
\midrule
\multicolumn{2}{l}{\textbf{Mean Absolute Error}} & \multicolumn{4}{c}{27.0\%} \\
\bottomrule
\multicolumn{6}{l}{\scriptsize All drugs classified as ``low confidence'' / extrapolation---expected and correct behavior.}
\end{tabular}
\end{table}

\end{document}
